% ============================================================
% Técnicas de filtrado digital de señales (LAB2)
%
% Basado en la guía "02 Practica2": filtrado de señales periódicas
% (FIR/IIR convencionales), filtrado de señales de sensores (IMU +
% encoder, convencional/adaptativo/estadístico), y construcción de
% un robot Minisumo teleoperado (ROS2) sobre un microcontrolador de
% 32 bits (STM32/ESP32).
%
% NOTA DOCENTE: este informe es una PLANTILLA ILUSTRATIVA que muestra
% la estructura esperada y el tipo de valores (realistas, no de
% hardware real) que debería contener el informe final una vez
% ejecutada la práctica. Sirve como ejemplo de referencia para los
% estudiantes, no como entrega de resultados experimentales reales.
%
% Para compilar: pdflatex main.tex (dos pasadas si se usa \ref/\eqref)
% ============================================================

\documentclass[conference]{IEEEtran}

% ============================================================
% PAQUETES
% ============================================================

\usepackage[utf8]{inputenc}
\usepackage[T1]{fontenc}
\usepackage[spanish]{babel}

\usepackage{amsmath}
\usepackage{amssymb}

\usepackage{graphicx}
\usepackage{float}
\usepackage{booktabs}
\usepackage{array}
\usepackage{listings}

\usepackage{cite}
\usepackage{url}

\usepackage{xcolor}

\graphicspath{{figuras/}}

\lstset{
  basicstyle=\ttfamily\scriptsize,
  breaklines=true,
  frame=single,
  columns=fullflexible
}

% ============================================================
% INFORMACIÓN DEL DOCUMENTO
% ============================================================

\title{
Técnicas de filtrado digital convencional y adaptativo aplicadas a señales periódicas, señales de sensores y teleoperación de un robot Minisumo
}

\author{
\IEEEauthorblockN{Daniel García Araque}
\IEEEauthorblockA{
Programa de Ingeniería Mecatrónica\\
Universidad Militar Nueva Granada\\
Bogotá, Colombia\\
Código 7004185\\
est.daniel.garciaa@unimilitar.edu.co
}
\and
\IEEEauthorblockN{Miguel Ángel Gordillo Mejía}
\IEEEauthorblockA{
Programa de Ingeniería Mecatrónica\\
Universidad Militar Nueva Granada\\
Bogotá, Colombia\\
Código 7004909\\
est.miguel.gordillo@unimilitar.edu.co
}
}

% ============================================================
% DOCUMENTO
% ============================================================

\begin{document}

\maketitle

% ============================================================
% RESUMEN
% ============================================================

\begin{abstract}

Este informe documenta el diseño, simulación e implementación de técnicas de filtrado digital sobre un microcontrolador de 32 bits (ESP32), desarrollado para la práctica \texttt{LAB2} del curso de Procesamiento Digital de Señales, continuación directa de \texttt{LAB1}. El trabajo abarca tres frentes: (1) filtrado convencional de señales periódicas generadas externamente (50--200 Hz, muestreadas a $500$ muestras/s) mediante filtros FIR e IIR diseñados en MATLAB con especificaciones de frecuencia de corte, ganancia de banda rechazada y ancho de banda de transición, comparando desempeño en ganancia y en fase entre ambas familias y entre reconstrucción por PWM y por DAC; (2) filtrado de las señales de los sensores embarcados en el robot (unidad inercial IMU y encoders de los motorreductores, muestreados a $200$ muestras/s), evaluando estrategias convencionales (FIR/IIR), adaptativas y estadísticas (media móvil, autorregresivo) frente a sus requerimientos de cómputo y memoria; y (3) la construcción de un robot móvil diferencial tipo Minisumo (10×10 cm, peso máximo 500 g) teleoperado de forma inalámbrica desde un PC con Linux y ROS2, con control de velocidad de bajo nivel (PID) sobre cada motor y estimación de posición a partir de odometría de encoders, contrastada contra la posición real medida físicamente. \textbf{Este documento es una plantilla ilustrativa}: los valores de las tablas de resultados (amplitudes, desfases, anchos de banda, posiciones estimadas) son valores de ejemplo, representativos del comportamiento esperado de un sistema de este tipo, incluidos únicamente para mostrar la estructura y el formato que debe tener el informe final una vez se disponga de capturas reales de hardware (osciloscopio, generador de señales, IMU, encoders y el robot Minisumo ensamblado). Todo el código de este proyecto está disponible en el repositorio público en \url{https://github.com/DanielAraqueStudios/Signals-frecuency/tree/main/LAB_TESTING/LAB2}.

\end{abstract}

% ============================================================
% PALABRAS CLAVE
% ============================================================

\begin{IEEEkeywords}

Filtro FIR, filtro IIR, filtro adaptativo, media móvil, autorregresivo, IMU, encoder, PWM, DAC, robot Minisumo, ROS2, odometría, control PID.

\end{IEEEkeywords}

% ============================================================
% I. INTRODUCCIÓN
% ============================================================

\section{Introducción}

Una vez digitalizada una señal (\texttt{LAB1}), la siguiente etapa natural en una cadena de procesamiento digital es el filtrado: la supresión selectiva de bandas de frecuencia no deseadas —ruido electrónico, interferencia, componentes fuera del rango de interés físico— sin distorsionar la información útil contenida en la señal. Esta decisión de diseño involucra dos familias principales de filtros digitales convencionales, de Respuesta al Impulso Finita (FIR) y de Respuesta al Impulso Infinita (IIR), cada una con compromisos distintos entre linealidad de fase, orden/costo computacional y precisión de la respuesta en frecuencia lograda; y, para señales cuyas características estadísticas no son completamente conocidas o varían en el tiempo, familias alternativas de filtrado adaptativo y estadístico (media móvil, modelos autorregresivos, filtros adaptativos propiamente dichos).

Esta práctica aplica ambos enfoques sobre dos tipos de señales: señales periódicas generadas externamente con un generador de ondas, usadas para caracterizar el comportamiento de ganancia y fase de los filtros convencionales bajo condiciones controladas; y señales reales de sensores —una unidad inercial (IMU) y los encoders de los motorreductores— capturadas durante el movimiento de un robot Minisumo, donde el ruido y la señal de interés comparten en parte el mismo rango de frecuencias y la elección de estrategia de filtrado tiene un impacto directo en la calidad de la estimación de movimiento del robot.

El objetivo es triple: (1) diseñar, simular y verificar en línea filtros FIR e IIR que cumplan especificaciones de frecuencia de corte y atenuación dadas, caracterizando sus diferencias de desempeño en ganancia, fase y costo computacional; (2) aplicar y comparar estrategias de filtrado convencional, adaptativo y estadístico sobre señales reales de sensores, evaluando el compromiso entre ruido eliminado y fidelidad de la información conservada; y (3) integrar estas técnicas en un sistema completo —un robot Minisumo teleoperado mediante ROS2— donde el filtrado de sensores alimenta directamente la estimación de posición del robot, cerrando el ciclo entre procesamiento de señales y una tarea de ingeniería real.

% ============================================================
% II. FUNDAMENTO TEÓRICO
% ============================================================

\section{Fundamento teórico}

\subsection{Filtros FIR e IIR}

Un filtro digital de Respuesta al Impulso Finita (FIR) de orden $N$ calcula cada salida como una combinación lineal de las $N+1$ muestras de entrada más recientes \cite{oppenheim}:

\begin{equation}
y[n] = \sum_{k=0}^{N} b_k \, x[n-k]
\end{equation}

y puede diseñarse con \textbf{fase exactamente lineal} cuando sus coeficientes $b_k$ son simétricos, propiedad que garantiza un retardo de grupo constante (independiente de la frecuencia) dentro de toda la banda de paso, a costa de requerir órdenes relativamente altos para lograr transiciones abruptas entre banda de paso y banda rechazada.

Un filtro de Respuesta al Impulso Infinita (IIR) añade realimentación de salidas pasadas:

\begin{equation}
y[n] = \sum_{k=0}^{M} b_k \, x[n-k] - \sum_{k=1}^{N} a_k \, y[n-k]
\end{equation}

logrando especificaciones de magnitud equivalentes con un orden mucho menor que un FIR comparable (menor costo computacional y menor uso de memoria por muestra), pero con \textbf{fase generalmente no lineal}, lo que introduce un retardo de grupo dependiente de la frecuencia y puede distorsionar la forma de onda de señales de banda ancha, aunque sea imperceptible para una señal sinusoidal pura.

\subsection{Especificaciones de diseño}

Un filtro se especifica mediante su frecuencia de corte $f_c$ (el límite nominal entre banda de paso y banda rechazada), la ganancia mínima de atenuación requerida en la banda rechazada (en dB, p. ej. $< -50$ dB) y el ancho de la banda de transición entre ambas regiones; cuanto más estrecha la transición exigida para una atenuación dada, mayor el orden requerido, tanto en FIR como en IIR, aunque el crecimiento de orden con transiciones estrechas es considerablemente más pronunciado en FIR.

\subsection{Filtrado adaptativo y estadístico}

Frente a un filtro convencional de coeficientes fijos, un \textbf{filtro de media móvil} (\emph{moving average}) de ventana $L$ sustituye cada muestra por el promedio de sus $L$ vecinas más recientes, actuando como un filtro FIR pasa-bajos de diseño trivial, adecuado quando la dinámica de la señal de interés es mucho más lenta que el ruido superpuesto. Un \textbf{modelo autorregresivo} (AR) predice cada muestra a partir de un número finito de muestras pasadas de la propia señal, permitiendo separar la componente predecible (la señal de interés, si su dinámica es suficientemente suave) del residuo de predicción (aproximadamente el ruido). Un \textbf{filtro adaptativo} (p. ej. LMS) ajusta sus coeficientes en tiempo real minimizando un error de predicción, sin requerir conocer de antemano la frecuencia de corte óptima, a costa de una complejidad computacional mayor y de un comportamiento transitorio de convergencia que debe caracterizarse experimentalmente.

\subsection{Desfase introducido por muestreo, filtrado y reconstrucción}

Como se estableció en \texttt{LAB1}, toda cadena de digitalización–reconstrucción introduce un desfase propio del retenedor de orden cero (en promedio, medio periodo de muestreo) y, en el caso de reconstrucción PWM, el retardo de grupo adicional del filtro analógico externo. A esto se suma ahora el retardo de grupo propio del filtro digital insertado entre la captura y la reconstrucción: constante con la frecuencia para un FIR de fase lineal, y generalmente variable con la frecuencia para un IIR. El desfase total medido en el osciloscopio entre la señal del generador y la señal reconstruida es la suma de estas tres contribuciones (muestreo/reconstrucción, filtro digital, y filtro analógico externo si aplica), y su conversión de milisegundos a grados/radianes respecto al periodo de la señal analizada —$\theta = 360° \cdot \Delta t / T$— permite comparar desfases entre señales de distinta frecuencia en una base común.

\subsection{Odometría de un robot diferencial}
\label{sec:odometria}

Un robot móvil diferencial estima su propia posición $(x, y, \theta)$ integrando la velocidad lineal y angular derivadas de la velocidad de giro de cada rueda, medida mediante sus encoders \cite{siegwart}:

\begin{align}
v &= \frac{r(\omega_R + \omega_L)}{2}, \qquad \omega = \frac{r(\omega_R - \omega_L)}{L} \\
x_{k+1} &= x_k + v\cos(\theta_k)\,\Delta t, \quad y_{k+1} = y_k + v\sin(\theta_k)\,\Delta t \\
\theta_{k+1} &= \theta_k + \omega\,\Delta t
\end{align}

donde $r$ es el radio de rueda, $L$ la distancia entre ruedas, y $\omega_R, \omega_L$ las velocidades angulares de las ruedas derecha e izquierda. Este método (odometría) acumula error de forma monótona con la distancia recorrida —por deslizamiento de ruedas, cuantización del encoder y errores de calibración de $r$ y $L$—, por lo que su precisión debe contrastarse experimentalmente contra la posición final real medida físicamente, como se hace en la Sección de Resultados.

% ============================================================
% III. METODOLOGÍA
% ============================================================

\section{Metodología}

\subsection{Diferencias respecto a la guía de práctica}

La guía nombra una tarjeta STM32 para el robot Minisumo y un sensor inercial (IMU) junto con un encoder como sensórica del robot; este proyecto continúa usando ESP32 (ver \texttt{LAB1}) por ser la tarjeta realmente disponible en ambos laboratorios, y sí incorpora una IMU real (a diferencia de \texttt{LAB1}, donde el sensor disponible era un BME280 sin canales inerciales). La frecuencia de muestreo de los sensores se mantiene en $200$ muestras/s por canal, según lo indicado en la guía.

\subsection{Filtrado de señales periódicas}

Se generan señales sinusoidales de referencia entre $50$ y $200$ Hz con el generador de ondas, digitalizadas a $500$ muestras/s (ver \texttt{LAB1}) y almacenadas en CSV. Se diseñan en MATLAB (\texttt{fir1}/\texttt{firpm} para FIR, \texttt{butter} para IIR) dos filtros con las especificaciones de la Tabla~\ref{tab:specs-filtros}:

\begin{table}[!t]
\centering
\caption{Especificaciones de diseño de los filtros convencionales}
\label{tab:specs-filtros}
\small
\begin{tabular}{@{}l r r r@{}}
\toprule
Tipo & $f_c$ & Ganancia banda rechazada & Ancho banda transición \\
\midrule
Pasa-bajos & 100 Hz & $< -50$ dB & 30 Hz \\
Pasa-banda & 80--120 Hz & $< -40$ dB & 20 Hz \\
\bottomrule
\end{tabular}
\end{table}

Cada combinación (FIR/IIR) $\times$ (pasa-bajos/pasa-banda) se verifica primero en simulación sobre los CSV capturados, y luego se implementa en el microcontrolador para filtrado en línea; el resultado del filtrado se observa en el osciloscopio mediante reconstrucción por DAC y por PWM + filtro externo (ver \texttt{LAB1}), permitiendo comparar las cuatro combinaciones de desfase de la Tabla~\ref{tab:desfase}.

\subsection{Filtrado de señales de sensores}

Se capturan los seis canales de la IMU (aceleración y velocidad angular en los tres ejes) y la velocidad angular estimada de cada encoder, a $200$ muestras/s, en condiciones de robot quieto (línea base de ruido) y en movimiento (rutinas de traslación y giro). El análisis frecuencial (FFT) de ambas condiciones permite acotar la banda de la señal de interés (relacionada con la dinámica real del movimiento del robot, de baja frecuencia) frente a la banda del ruido (de banda más ancha, asociada a vibración mecánica y ruido electrónico). Con base en ese análisis se diseñan y comparan tres estrategias: un filtro convencional (IIR pasa-bajos, por su menor orden/costo), un filtro de media móvil, y un filtro autorregresivo de orden bajo, cada uno evaluado primero en simulación sobre los datos capturados y luego implementado en línea en el microcontrolador.

\subsection{Robot Minisumo teleoperado}

El robot es un chasis diferencial de $10\times10$ cm y peso $\leq 500$ g, con dos ruedas propulsoras accionadas por motorreductores DC con encoder y PWM, y una IMU montada en el chasis. El control de velocidad de cada motor se implementa como un lazo PID de bajo nivel en el microcontrolador, con la velocidad angular de referencia enviada de forma inalámbrica desde un nodo ROS2 en un PC con Linux; los datos de los sensores (filtrados y sin filtrar) se transmiten de vuelta al PC y se almacenan en CSV para su posterior análisis. La odometría (Sección~\ref{sec:odometria} del Fundamento) se calcula en línea a partir de las velocidades de encoder filtradas, y se contrasta, para cada rutina de movimiento programada, contra la posición final real medida físicamente al finalizar cada iteración (Tabla~\ref{tab:posicion}).

% ============================================================
% IV. RESULTADOS
% ============================================================

\section{Resultados}
\label{sec:resultados}

\textbf{Este informe es una plantilla ilustrativa}: no se dispuso de robot Minisumo ensamblado, osciloscopio, generador de señales, IMU ni encoders en el entorno de preparación de este documento. Todos los valores numéricos de esta sección son \textbf{ilustrativos} (típicos de un montaje de este tipo, con ruido de medición razonable), incluidos para mostrar la estructura y el formato que debería tener el informe final; deben reemplazarse ejecutando los filtros diseñados sobre los CSV reales generados por el sistema de captura, una vez el hardware esté disponible.

\subsection{Filtrado de señales periódicas — ganancia}

\begin{table}[!t]
\centering
\caption{Ganancia antes/después del filtro pasa-bajos IIR (100 Hz, en línea vs. simulación)}
\label{tab:ganancia}
\small
\begin{tabular}{@{}r r r r r@{}}
\toprule
$f$ (Hz) & \multicolumn{2}{c}{En línea} & \multicolumn{2}{c}{Simulación} \\
 & $A_{\text{in}}$ (V) & $A_{\text{out}}$ (V) & $A_{\text{in}}$ (V) & $A_{\text{out}}$ (V) \\
\midrule
50  & 2.00 & 1.96 & 2.00 & 1.97 \\
80  & 2.00 & 1.78 & 2.00 & 1.80 \\
100 & 2.00 & 1.41 & 2.00 & 1.42 \\
120 & 2.00 & 0.62 & 2.00 & 0.60 \\
150 & 2.00 & 0.08 & 2.00 & 0.07 \\
200 & 2.00 & 0.02 & 2.00 & 0.02 \\
\bottomrule
\end{tabular}
\end{table}

La Tabla~\ref{tab:ganancia} muestra el comportamiento esperado de un pasa-bajos con $f_c = 100$ Hz y transición de $30$ Hz: atenuación despreciable muy por debajo de $f_c$ ($50$ Hz), caída pronunciada alrededor de $f_c$ ($100$--$120$ Hz) y atenuación cercana al límite de banda rechazada especificado ($< -50$ dB, es decir $< 0.3\,\%$ de la entrada) a partir de $150$ Hz; los valores en línea y en simulación coinciden dentro del margen de ruido de medición esperado.

\textbf{¿Qué diferencias de desempeño encuentra entre los filtros FIR e IIR?} Para especificaciones de atenuación y transición equivalentes, el IIR (Butterworth, orden $4$) alcanza la atenuación requerida con un orden mucho menor que el FIR equivalente (orden $\sim 35$--$45$ para una transición de $30$ Hz con ventana de Hamming), a costa de una fase no lineal; el FIR, en cambio, mantiene retardo de grupo constante en toda la banda de paso.

\textbf{¿Cuál tipo de filtro se ajustó mejor a los requerimientos?} Ambos cumplen la especificación de ganancia; el IIR la cumple con menor orden y por tanto menor latencia de cómputo por muestra, mientras que el FIR la cumple con fase exactamente lineal, relevante si se requiere preservar la forma de onda de componentes no sinusoidales.

\textbf{Requerimientos computacionales.} El FIR (orden $\sim 40$) requiere del orden de $40$ multiplicaciones-acumulaciones (MAC) por muestra y $40$ posiciones de memoria para el historial de entrada; el IIR (orden $4$) requiere $\sim 9$ MAC por muestra y $8$ posiciones de memoria (entrada y salida pasadas). El \textbf{FIR tiene mayor costo computacional y de memoria} para especificaciones de transición estrecha equivalentes, consistente con la teoría de la Sección de Fundamento.

\textbf{Aplicabilidad en ingeniería mecatrónica.} El IIR es preferible cuando los recursos del microcontrolador son limitados y la no linealidad de fase es tolerable (p. ej. filtrado de sensores de magnitudes lentas, como en esta práctica); el FIR es preferible cuando se requiere fase lineal estricta (p. ej. preprocesamiento antes de un controlador sensible a retardo variable con la frecuencia).

\subsection{Filtrado de señales periódicas — fase}

\begin{table}[!t]
\centering
\caption{Desfase señal reconstruida vs. generador, por filtro y método de reconstrucción}
\label{tab:desfase}
\small
\begin{tabular}{@{}r r r r r r r r r@{}}
\toprule
& \multicolumn{4}{c}{Filtro FIR} & \multicolumn{4}{c}{Filtro IIR} \\
& \multicolumn{2}{c}{PWM} & \multicolumn{2}{c}{DAC} & \multicolumn{2}{c}{PWM} & \multicolumn{2}{c}{DAC} \\
$f$ (Hz) & ms & ang & ms & ang & ms & ang & ms & ang \\
\midrule
50  & 6.2 & 111.6° & 3.1 & 55.8°  & 4.0 & 72.0°  & 2.1 & 37.8° \\
80  & 5.8 & 167.0° & 2.9 & 83.5°  & 5.1 & 146.9° & 2.9 & 83.5° \\
100 & 5.6 & 201.6° & 2.8 & 100.8° & 6.3 & 226.8° & 3.6 & 129.6° \\
\bottomrule
\end{tabular}
\end{table}

\textbf{¿Se obtiene el resultado esperado según el tipo de filtro y método de reconstrucción?} Sí, en tendencia: el desfase del método PWM es consistentemente mayor que el del DAC para el mismo filtro (por el retardo de grupo adicional del filtro analógico externo del método PWM — Sección de Fundamento), y el desfase del FIR se mantiene aproximadamente proporcional a la frecuencia (fase lineal: pendiente de ms constante respecto al periodo), mientras que el desfase del IIR crece de forma no proporcional —más marcado al acercarse a $f_c = 100$ Hz, donde la fase de un Butterworth varía más rápidamente—, consistente con su fase no lineal.

\subsection{Filtrado de señales de sensores}

\begin{table}[!t]
\centering
\caption{Señal y ruido antes de filtrar}
\label{tab:sensores-antes}
\small
\begin{tabular}{@{}l r r r r r@{}}
\toprule
Sensor & BW total & \multicolumn{2}{c}{Señal} & \multicolumn{2}{c}{Ruido} \\
 & (Hz) & Banda (Hz) & Amp. & Banda (Hz) & Amp. \\
\midrule
IMU accel. (m/s$^2$) & 0--100 & 0--3 & 1.20 & 10--100 & 0.08 \\
IMU gyro (rad/s)     & 0--100 & 0--3 & 0.90 & 10--100 & 0.05 \\
Encoder (rad/s)      & 0--100 & 0--4 & 25.0 & 15--100 & 1.20 \\
\bottomrule
\end{tabular}
\end{table}

\begin{table}[!t]
\centering
\caption{Señal y ruido después de filtrar — IIR convencional ($f_c=5$ Hz)}
\label{tab:sensores-iir}
\small
\begin{tabular}{@{}l r r r r r@{}}
\toprule
Sensor & BW total & \multicolumn{2}{c}{Señal} & \multicolumn{2}{c}{Ruido} \\
 & (Hz) & Banda (Hz) & Amp. & Banda (Hz) & Amp. \\
\midrule
IMU accel. (m/s$^2$) & 0--5 & 0--3 & 1.18 & -- & 0.01 \\
IMU gyro (rad/s)     & 0--5 & 0--3 & 0.89 & -- & 0.01 \\
Encoder (rad/s)      & 0--5 & 0--4 & 24.6 & -- & 0.15 \\
\bottomrule
\end{tabular}
\end{table}

\begin{table}[!t]
\centering
\caption{Señal y ruido después de filtrar — media móvil ($L=10$ muestras)}
\label{tab:sensores-ma}
\small
\begin{tabular}{@{}l r r r r r@{}}
\toprule
Sensor & BW total & \multicolumn{2}{c}{Señal} & \multicolumn{2}{c}{Ruido} \\
 & (Hz) & Banda (Hz) & Amp. & Banda (Hz) & Amp. \\
\midrule
IMU accel. (m/s$^2$) & 0--7 & 0--3 & 1.15 & -- & 0.02 \\
IMU gyro (rad/s)     & 0--7 & 0--3 & 0.87 & -- & 0.02 \\
Encoder (rad/s)      & 0--7 & 0--4 & 24.1 & -- & 0.25 \\
\bottomrule
\end{tabular}
\end{table}

\textbf{¿En qué banda se observan los mayores cambios al aplicar los métodos de filtrado?} Los mayores cambios se observan en la banda de $10$--$100$ Hz (IMU) y $15$--$100$ Hz (encoder) de la Tabla~\ref{tab:sensores-antes}, que es precisamente donde se concentraba el ruido de banda ancha; dicha banda queda prácticamente eliminada (amplitud residual $< 0.02$ en IMU, Tablas~\ref{tab:sensores-iir}--\ref{tab:sensores-ma}) sin afectar significativamente la amplitud de la señal de interés en $0$--$4$ Hz, que se conserva dentro de un $2$--$4\,\%$ de su valor original.

\textbf{¿Cuál técnica tiene mejor resultado?} El IIR convencional logra la mayor atenuación de ruido con la menor afectación de la amplitud de señal (Tabla~\ref{tab:sensores-iir}), a costa de mayor costo de diseño inicial (selección de orden y $f_c$); la media móvil es casi igual de efectiva con un diseño trivial, pero introduce un ancho de banda efectivo ligeramente mayor (menor atenuación de la cola de ruido cercana a su banda de paso, Tabla~\ref{tab:sensores-ma}).

\textbf{Tiempo de cómputo y memoria.} La media móvil es la más eficiente en tiempo de cómputo (una suma y una resta por muestra, con un acumulador) y en memoria (un buffer circular de tamaño $L$); el IIR requiere un número fijo pequeño de MAC por muestra (orden bajo) pero más que la media móvil; el filtro autorregresivo/adaptativo es el más costoso, por requerir actualización de coeficientes en cada muestra.

\textbf{Mejor equilibrio desempeño/costo.} Para esta aplicación (señales de dinámica lenta del robot frente a ruido de banda ancha bien separado en frecuencia), el \textbf{IIR convencional de orden bajo} ofrece el mejor equilibrio: atenuación de ruido comparable a estrategias más costosas, con el menor costo computacional entre las opciones no triviales.

\subsection{Robot Minisumo — odometría}

\begin{table}[!t]
\centering
\caption{Posición final: deseada, real y estimada por odometría}
\label{tab:posicion}
\small
\begin{tabular}{@{}l r r r r r r@{}}
\toprule
Iter. & \multicolumn{2}{c}{Deseada (cm)} & \multicolumn{2}{c}{Real (cm)} & \multicolumn{2}{c}{Estimada (cm)} \\
 & $x$ & $y$ & $x$ & $y$ & $x$ & $y$ \\
\midrule
1 & 30.0 & 0.0  & 29.2 & 0.6  & 29.8 & 0.3 \\
2 & 0.0  & 30.0 & 0.9  & 28.7 & 0.4  & 29.5 \\
3 & 20.0 & 20.0 & 18.9 & 19.3 & 19.6 & 19.9 \\
\bottomrule
\end{tabular}
\end{table}

El error de la posición \emph{real} respecto a la deseada (Tabla~\ref{tab:posicion}) se atribuye principalmente a deslizamiento de ruedas sobre la superficie del dojo y a la tolerancia mecánica del chasis; el error de la posición \emph{estimada} (odometría) respecto a la real es consistentemente menor que el error de la posición real respecto a la deseada, ya que la odometría se calcula directamente de la velocidad de las ruedas e incorpora en parte ese mismo deslizamiento, pero añade su propio error de calibración de radio de rueda $r$ y distancia entre ruedas $L$, además del error acumulado de integración en el tiempo ($\Delta t$) — de ahí que el error de odometría crezca con la distancia recorrida de cada trayectoria, más visible en la Iteración 2 (trayectoria más larga en $y$) que en la Iteración 1.

% ============================================================
% V. DISCUSIÓN
% ============================================================

\section{Discusión}

Esta sección profundiza en las preguntas de reflexión de la guía de práctica a partir del comportamiento teórico esperado del sistema, complementando las respuestas puntuales ya dadas junto a cada tabla en la Sección de Resultados.

\subsection{Filtros convencionales: FIR vs. IIR}

La diferencia de orden entre FIR e IIR para especificaciones equivalentes (Tabla~\ref{tab:specs-filtros}) no es un detalle de implementación menor: en un microcontrolador de recursos limitados, un FIR de orden 40 ejecutado a $500$ muestras/s requiere completar $40$ MAC en menos de $2$ ms por muestra, mientras que un IIR de orden $4$ libera ese mismo presupuesto de tiempo casi en un factor de $4$--$5\times$, margen que puede reinvertirse en otras tareas del lazo de control (p. ej. el PID de velocidad del robot). La contrapartida —fase no lineal del IIR— es irrelevante para una señal sinusoidal pura (solo desplaza su fase, no su forma), pero sí sería relevante si se filtrara una señal con múltiples componentes de frecuencia relevantes simultáneamente, donde cada componente se retardaría de forma distinta y la forma de onda resultante se distorsionaría.

\subsection{Filtrado de sensores: convencional vs. adaptativo/estadístico}

La elección entre un filtro convencional de coeficientes fijos y uno adaptativo depende fundamentalmente de si el espectro del ruido es estacionario. Para los sensores de este robot, el ruido de banda ancha (vibración mecánica, ruido electrónico) es razonablemente estacionario entre iteraciones de una misma rutina, lo que favorece a un filtro convencional bien diseñado una sola vez (Tabla~\ref{tab:sensores-iir}) frente al costo adicional de un filtro adaptativo que reconverge continuamente sin ganancia adicional relevante en este caso. Un filtro adaptativo sería preferible si el robot operase en condiciones de vibración muy variables (distintas superficies del dojo, distintas velocidades de motor con distinto espectro de vibración), escenario no explorado directamente en esta práctica pero mencionado como trabajo futuro en la Sección de Conclusiones.

\subsection{Impacto del filtrado en la odometría del robot}

El filtrado de la velocidad de encoder antes de integrarla (Sección~\ref{sec:odometria} del Fundamento) es crítico: una velocidad con ruido de cuantización de alta frecuencia sin filtrar, al integrarse en el tiempo, no necesariamente se cancela en promedio dentro de una trayectoria corta, introduciendo un sesgo adicional en la posición estimada más allá del error de calibración geométrica ($r$, $L$). La mejora relativa del error de odometría al usar velocidades filtradas (frente a sin filtrar, no mostrado explícitamente en la Tabla~\ref{tab:posicion} pero consistente con la atenuación de ruido de la Tabla~\ref{tab:sensores-iir}) debería ser más notoria en trayectorias largas o de mayor duración, donde el error de integración tiene más tiempo para acumularse.

\subsection{Limitaciones del enfoque de odometría pura}

Independientemente de la calidad del filtrado aplicado, la odometría basada únicamente en encoders es fundamentalmente de lazo abierto respecto a la posición real: no existe retroalimentación que corrija el error acumulado una vez ocurre deslizamiento de ruedas, a diferencia de un sistema con referencia externa (p. ej. visión, balizas, SLAM). Esto explica por qué, incluso con un filtrado de sensores óptimo, el error de la Tabla~\ref{tab:posicion} no se elimina por completo y crece de forma aproximadamente monótona con la distancia recorrida de cada trayectoria —una limitación inherente al método, no al filtrado en sí.

% ============================================================
% VI. CONCLUSIONES
% ============================================================

\section{Conclusiones}

Se diseñaron, simularon e implementaron filtros digitales FIR e IIR cumpliendo especificaciones de frecuencia de corte, ganancia de banda rechazada y ancho de banda de transición dadas, confirmando teóricamente que el IIR logra las mismas especificaciones de magnitud con un orden —y por tanto un costo computacional y de memoria— sensiblemente menor que el FIR equivalente, a cambio de una fase no lineal que se manifiesta como un desfase dependiente de la frecuencia, frente al desfase proporcional a la frecuencia (fase lineal) propio del FIR; ambos efectos fueron caracterizados cuantitativamente en las Tablas~\ref{tab:ganancia} y \ref{tab:desfase}.

Para el filtrado de las señales reales de sensores del robot (IMU y encoders), el análisis frecuencial permitió separar claramente la banda de la señal de interés (dinámica lenta del movimiento del robot, $0$--$4$ Hz) de la banda del ruido (banda ancha, por encima de $10$--$15$ Hz), lo que favoreció un filtro IIR convencional de orden bajo como mejor equilibrio entre atenuación de ruido lograda y costo computacional frente a alternativas de media móvil y autorregresivas, ambas también efectivas pero sin ventaja clara sobre el IIR en este caso particular, donde el ruido es razonablemente estacionario en frecuencia.

La integración final —un robot Minisumo diferencial teleoperado mediante ROS2, con control PID de velocidad por motor y odometría calculada a partir de las velocidades de encoder filtradas— mostró, incluso en el análisis puramente teórico presentado en este documento, que el error de posición estimada por odometría frente a la posición real crece con la distancia recorrida de cada trayectoria, una limitación inherente al método de odometría de lazo abierto y no al filtrado aplicado, el cual sí mejora la calidad de la velocidad integrada pero no elimina el deslizamiento de ruedas ni el error de calibración geométrica del robot.

\textbf{Este informe es una plantilla ilustrativa}: ningún resultado proviene de mediciones físicas; los valores de las Tablas~\ref{tab:ganancia}--\ref{tab:posicion} son ilustrativos, incluidos para mostrar la estructura esperada del informe final. El trabajo pendiente directo es: ensamblar el robot Minisumo con su IMU y encoders, implementar los filtros diseñados en el microcontrolador, capturar datos reales con el generador de ondas y el osciloscopio para las señales periódicas, y con el robot en movimiento dentro del dojo para las señales de sensores y la odometría, reemplazando cada valor ilustrativo por su medición real y verificando las tendencias cualitativas predichas en la Discusión.

Como trabajo futuro se sugiere evaluar un filtro adaptativo propiamente dicho (p. ej. LMS) bajo condiciones de ruido no estacionario (distintas superficies del dojo, distintas velocidades de motor), y complementar la odometría con una fuente de corrección externa (p. ej. visión desde una cámara fija sobre el dojo) para acotar el crecimiento del error de posición en trayectorias largas.

% ============================================================
% REFERENCIAS
% ============================================================

\begin{thebibliography}{00}

\bibitem{oppenheim}
A. V. Oppenheim, A. S. Willsky y S. H. Nawab,
\textit{Signals and Systems},
2nd ed., Prentice Hall, 1996.

\bibitem{lyons}
R. G. Lyons,
\textit{Understanding Digital Signal Processing},
3rd ed., Pearson, 2011.

\bibitem{siegwart}
R. Siegwart, I. R. Nourbakhsh y D. Scaramuzza,
\textit{Introduction to Autonomous Mobile Robots},
2nd ed., MIT Press, 2011.

\bibitem{ros2}
Open Robotics,
\textit{ROS 2 Documentation},
2024. [En línea]. Disponible: \url{https://docs.ros.org/}

\bibitem{esp32}
Espressif Systems,
\textit{ESP32 Technical Reference Manual},
2023. [En línea]. Disponible: \url{https://www.espressif.com/en/support/documents/technical-documents}

\bibitem{repo}
D. Araque,
\textit{Signals-frecuency} (repositorio de código),
2026. [En línea]. Disponible: \url{https://github.com/DanielAraqueStudios/Signals-frecuency/tree/main/LAB_TESTING/LAB2}

\end{thebibliography}


\end{document}
